\documentclass[11pt]{article}

\usepackage[margin=1in]{geometry}
\usepackage{times}
\usepackage[T1]{fontenc}
\usepackage[utf8]{inputenc}
\usepackage{microtype}
\usepackage{setspace}
\usepackage{parskip}

\setstretch{1.10}
\setlength{\parindent}{0pt}
\setlength{\parskip}{6pt plus 1pt}

\pagenumbering{gobble}

\begin{document}

\begin{center}
  {\Large \textbf{Research Statement}} \\[3pt]
  \textit{Xueyang Liu \,---\, ESADE Business School (Universitat Ramon Llull)} \\
  \textit{Application: 1st Zurich Summer School in AI \& Applied Economics}
\end{center}

\vspace{4pt}

I am a first-year PhD student in Management Sciences at ESADE Business School
(Universitat Ramon Llull, Barcelona), supervised by Dr.\ Emanuele Rizzo. My
research uses novel, AI-era sources of firm-level data to study how firms'
web-level configurations shape investor information acquisition and capital
market outcomes. The project I would bring to the summer school examines
whether firms whose public websites restrict AI search crawlers become
systematically less visible in AI-generated answers, and whether this reduced
visibility translates into measurable changes in the firm's information
environment.

The project rests on an original dataset I have constructed from the technical
configuration of corporate websites. For each of approximately 9{,}500 US
public firms, I parse the website's \texttt{robots.txt} file according to
RFC~9309 and match each declared sitemap URL against the firm's permissions
for sixteen AI search crawlers, including OpenAI's OAI-SearchBot, Anthropic's
Claude-SearchBot, and PerplexityBot. The resulting measure, sitemap--robots
overlap, identifies firms whose own declarations of crawl-discoverable content
conflict with their own access rules. This is a form of information friction
that does not appear in standard disclosure data and that firms may inherit
unintentionally from default website templates. I am now extending this
cross-section into a monthly panel for 2021--2025 using the Wayback Machine,
which will allow for a staggered difference-in-differences design around
firm-level policy changes.

To validate the mechanism, I am running a controlled experimental audit using
the OpenAI, Perplexity, and Anthropic APIs. For a matched-pair sample of
treated and control firms, I programmatically issue identical natural-language
queries and record whether each firm's official website appears among the
inline citations. Preliminary results from the ChatGPT and Claude pilots show
a consistent reduction of eight to eleven percentage points in the probability
of official-site citation for treated firms. I implemented the full pipeline
in Python, including the RFC~9309-compliant \texttt{robots.txt} parser, a
sitemap crawler with gzip and XML-index support, a Wayback Machine CDX
client, and the automated LLM-API harness that runs the citation audit.

The project fits the summer school's two core domains closely. On Topic~1
(empirical research with AI tools), the sitemap--robots overlap is a
measurement and dataset-construction exercise that applies an AI-era
concept---crawler-level access rules---to the classic corporate-finance
question of how information reaches investors. On Topic~2 (experiments and
AI), the API-based citation audit is a behavioral experiment on LLM agents,
in which the treatment varies whether the firm's content is accessible to the
agent, and the outcome measures what the agent produces in response to user
queries. The school's emphasis on identification, measurement, and inference
is directly useful to both components of the project.

I hope to use the week to sharpen the identification strategy for the panel
analysis, learn how other applied economists are handling LLM-agent
experiments, and receive feedback on the mapping from AI search visibility to
downstream capital-market outcomes. I am also eager to discuss with other
participants how AI-era measurement can be extended to other settings in
which firms' digital footprints shape market outcomes.

\end{document}
